For the horizon pair, a state vector that is invariant up to phase under the boost and cyclic and separating for the pair factor has the form $\cos\theta\,\ket{00} + \sin\theta\,\ket{11}$ with $\cos\theta\sin\theta \neq 0$ (the cyclic and separating hypothesis excludes the product vectors $\ket{01}$, $\ket{10}$, which are also invariant up to phase); its one-sided marginal is $\operatorname{diag}(\cos^{2}\theta, \sin^{2}\theta)$, the Gibbs state for the number operator at $\tan^{2}\theta = e^{-\beta\omega}$; the sample point $e^{-\omega/T} = \tfrac{9}{16}$ has $(\cos\theta, \sin\theta) = (\tfrac45, \tfrac35)$. Every automorphism of the mode factor intertwining the boost's phase flow $\operatorname{Ad}(e^{JtK})$, with $J = e_{2}e_{3}$ the complex unit of the radial atom (\claimref{synthesis.transverse-bivector-is-the-complex-unit-of-the-radial-atom}), acts on $\beta$ through $\{\pm 1\}$ only. The surface gravity is $\kappa = 1/4M$, the Euclidean section is smooth at the bifurcation sphere exactly for the period $2\pi/\kappa$, and (R, \claimref{singularity.r-regularity-fixed-locus}) imposes KMS at that period mode by mode. The supplied class (\claimref{singularity.l-boundary-mode-locality}) is ultralocal, and the s-wave array of Step 4e has weights $\omega_{k} = \kappa k$ with its sea level at $k = 0$ (by \claimref{singularity.q-joint-quasi-freeness-s-wave} in the source paper's convention, version 2 of premise (Q); version 3 adopts $\kappa(k + \tfrac12)$). \emph{Provenance:} \cite[Lem.~7.2, Thm.~7.3, Thm.~7.16, Prop.~7.18, Prem.~7.19, Prop.~7.25, Rem.~7.26, Prem.~7.30, Rem.~7.33, Prem.~8.1, Thm.~C.8, Rem.~C.9, Prop.~A.2]{ButlerT5} (\claimref{singularity.rate-rigidity-temperature-not-algebra-intrinsic}, \claimref{singularity.modular-descent}, \claimref{singularity.schmidt-forcing-proved-outright}, \claimref{singularity.euclidean-section-period}).

This restates \claimref{singularity.class-relative-no-go-no-supplied}, \claimref{singularity.computed}.
